\section{A control with a one-step modular potential}\label{sec:controls}

The freedom to choose the numerical control is decisive.  We use the
finite hierarchy
\[
 F_0(n)=2^n,\qquad F_{j+1}(n)=F_j^{\,n+1}(n),\qquad
 G(n)=F_5(n+100).
\]
The superscript on a function denotes iteration.  Put
\(p(n)=\lfloor n/2\rfloor\) and
\(L(n)=\operatorname{lcm}(1,\ldots,n)\), with \(L(0)=1\).
For a prescribed integer \(b\ge0\), define
\begin{align}
 f_*(n)&=G(n)!+2p(n)-2,\label{eq:factorial-control}\\
 f_b(n)&=\begin{cases}
 G(n),&n\le b,\\
 G(n)+((2p(n)-2-G(n))\bmod 2L(n)),&n>b.
 \end{cases}\label{eq:near-control}
\end{align}
All remainders are least nonnegative remainders.  The first definition
makes the divisibility transparent.  The second preserves an arbitrary
finite prefix of the published example and satisfies
\begin{equation}\label{eq:control-distance}
 0\le f_b(n)-G(n)<2L(n)\le2n!\quad(n>b).
\end{equation}
Below, \(f\) is either of these controls, and
\(B_f=0\) for \(f_*\), \(B_f=b+1\) for \(f_b\).

\begin{lemma}[Growth and modular decrement]\label{lem:control}
The functions \(f_*\) and \(f_b\) are even-valued primitive recursive
functions, majorize \(G\), and satisfy
\begin{equation}\label{eq:growth}
 f(n+1)\ge2^{2^{f(n)}},\qquad f(n)<G(n+1).
\end{equation}
They are eventually bounded by one fixed finite level of the hierarchy.
For every \(m\ge1\),
\begin{equation}\label{eq:potential}
 p(f(n))\equiv p(n)-1\pmod m
 \qquad(n\ge\max\{m,B_f\}).
\end{equation}
\end{lemma}
\begin{proof}
Factorial, least common multiple, and bounded remainder are primitive
recursive.  Every value of \(G\) is even.  The defining congruence for
\(f_b\) is even, so \(f_b\) is even too.  Write \(A=G(n)\).
Monotonicity of the iterates defining \(F_5\), including the extra iterate
at the larger argument, gives
\[
 G(n+1)\ge F_4(G(n)+1).
\]
For \(A\ge16\),
\[
 2A!\le2A^A<2^{A^2}<2^{2^A}.
\]
The fixed shift by 100 ensures this range at every input.  Also
\(G(n)\ge2n!\) and \(G(n)\ge2n+2\).  Consequently
\[
 G(n)\le f_b(n)<2G(n),\qquad
 G(n)\le f_*(n)\le2G(n)!.
\]
The latter lower bound holds also at \(n=0,1\), where subtracting two
from \(G(n)!\) still leaves a number above \(G(n)\).
The function \(F_4(A+1)\) dominates four successive exponentiations of
two beginning at \(A\).  The preceding estimates therefore imply
\eqref{eq:growth}, including the transition from the preserved prefix
to the corrected tail in \(f_b\).  For all sufficiently large \(n\),
\(G(n+1)=F_5(n+101)\le F_5(F_5(n))\le F_6(n)\), which is the
required eventual finite-level bound.

For \(f_*\) one has
\[
 p(f_*(n))=G(n)!/2+p(n)-1.
\]
If \(m\le n\), the integer \(2m\) is at most \(G(n)\), and hence
\(m\mid G(n)!/2\).  For \(f_b\), its definition gives
\(f_b(n)\equiv2p(n)-2\pmod{2L(n)}\) on the corrected tail.
Divide this congruence by two and use \(m\mid L(n)\).  This proves
\eqref{eq:potential}.
\end{proof}

The estimate \eqref{eq:growth} supplies the norm and separation margins
used in the canonical assignment.  Equation \eqref{eq:potential} has a
different purpose: it records the number of descent events without
computing that number.  Since the norm orbit is strictly increasing,
ordinary induction gives
\begin{equation}\label{eq:orbit-potential}
 p(f^i(t))\equiv p(t)-i\pmod m
 \quad\text{for every }i\ge0,
 \quad t\ge\max\{m,B_f\}.
\end{equation}
There is no bound on \(i\) in this identity.

\section{Complete descent counts and entry residues}\label{sec:residues}

For a nonzero limit \(\lambda<\varepsilon_0\), define
\[
 J_\lambda(t)=\psi_f(\lambda+t-N(\lambda)),
 \qquad t\ge N(\lambda).
\]
The greedy evaluation starts at \(\lambda+t-N(\lambda)\).  At a
nonfinite state \(h+k\), where \(h\) is its nonzero limit part, it
moves to \(M(h,f(N(h+k)))\).  Let \(D_\lambda(t)\) be the number
of such moves and \(Z_\lambda(t)\) the first finite state.  Thus the
finite predecessor chain after that state is not included in \(D\).
The maximal-predecessor and single-path lemmas give
\begin{equation}\label{eq:terminal-decomposition}
 u_i=f^i(t),\qquad Z_\lambda(t)=f^{D_\lambda(t)}(t),\qquad
 J_\lambda(t)=Z_\lambda(t)+D_\lambda(t).
\end{equation}
The use of the limit part in the predecessor call is essential.

\begin{theorem}[Uniform complete-count residues]\label{thm:counts}
For every \(m\ge1\) and every nonzero limit \(\lambda\), put
\(T=\max\{N(\lambda),m,B_f\}\).  For all \(t\ge T\),
\begin{align*}
 \lambda\ge\omega^2:\quad&
 D_\lambda(t)\equiv p(t)\pmod m,\quad
 Z_\lambda(t)\equiv0\pmod{2m},\quad
 J_\lambda(t)\equiv p(t)\pmod m;\\
 \lambda=\omega k:\quad&
 D_\lambda(t)=k,\quad
 Z_\lambda(t)\equiv2p(t)-2k\pmod{2m},\quad
 J_\lambda(t)\equiv2p(t)-k\pmod m.
\end{align*}
In particular, the first line gives a period \(2m\) and the explicit
table \(r\mapsto\lfloor r/2\rfloor\bmod m\) for every head
\(\lambda\ge\omega^2\).
\end{theorem}
\begin{proof}
Suppose \(\lambda\ge\omega^2\).  As long as the current limit part
is above \(\omega^2\), the ordinal \(\omega^2\), of norm three,
is an affordable predecessor.  The next limit part cannot fall below it.
Strict descent of limit parts therefore forces a visit to
\(\omega^2\).  Suppose this visit occurs at step \(i\), with norm
\(v=f^i(t)\).  Evenness of the next budget gives the exact predecessor
\[
 M(\omega^2,f(v))=\omega\cdot f(v)/2.
\]
An \(\omega\)-coefficient decreases by exactly one at each subsequent
nonfinite move, regardless of the finite tail.  Hence
\begin{equation}\label{eq:funnel-count}
 D_\lambda(t)=i+1+f(v)/2=i+1+p(f(v)).
\end{equation}
Equation \eqref{eq:orbit-potential} cancels the unknown \(i\), giving
\(D_\lambda(t)\equiv p(t)\).  Apply that equation once more at the
terminal event.  Since the terminal state is even,
\(p(Z_\lambda(t))\equiv0\pmod m\) implies
\(Z_\lambda(t)\equiv0\pmod{2m}\).  Now use
\eqref{eq:terminal-decomposition}.

For \(\lambda=\omega k\), exactly \(k\) nonfinite moves are made.
The same orbit identity and terminal decomposition give the second line.
These two cases exhaust the nonzero limits below \(\varepsilon_0\).
\end{proof}

The theorem uses one final passage through \(\omega^2\), rather than
an induction computing the count at each higher ordinal.  The actual
count can be enormous; its residue depends only on the original input.

For \(\lambda>\nu\), with \(\lambda\) a nonzero limit and
\(\nu\) a limit or zero, write \(E_{\lambda,\nu}(t)\) for the norm
at the first entry into \([\nu,\nu+\omega)\), and
\(\tau_{\lambda,\nu}(t)\) for its event count.  These functions
are defined only when that entry occurs.  On an affordable tail,
\begin{equation}\label{eq:entry-orbit}
 E_{\lambda,\nu}(t)=f^{\tau_{\lambda,\nu}(t)}(t).
\end{equation}
The case \(\nu=0\) is the terminal pair \((Z_\lambda,D_\lambda)\).
An identity interval is a separate convention with \(E(t)=t\) and
\(\tau(t)=0\).

\begin{proposition}[Entry tables]\label{prop:entry-residues}
Put \(c_\nu=\lceil N(\nu)/2\rceil\).  On a computable tail,
an interval of the form \(\lambda=\nu+\omega k\), \(k\ge1\),
has \(\tau(t)=k\) and \(E(t)=f^k(t)\).  Every other proper interval
has
\begin{equation}\label{eq:long-entry-residues}
 \tau(t)\equiv p(t)-c_\nu\pmod m,\qquad
 E(t)\equiv2c_\nu\pmod{2m}.
\end{equation}
A sufficient tail for these formulas is
\(t\ge\max\{N(\lambda),N(\nu)+3,m,B_f\}\).
\end{proposition}
\begin{proof}
Write \(\nu=P+\omega j\), where every exponent in \(P\), if any,
is at least two.  A proper interval is either the displayed finite
\(\omega\)-tail or has \(\lambda\ge P+\omega^2\).
In the latter case require
\[
 f(t)\ge N(P)+\max\{3,2j\}.
\]
This covers the norm bounds of the intermediate limit \(P+\omega^2\)
and the true target.
The stated tail implies this condition.  If the funnel is reached after
\(i\) events at norm \(v=f^i(t)\), its next \(\omega\)-coefficient
is \(f(v)/2-\lceil N(P)/2\rceil\).  Therefore
\[
 \tau(t)=i+1+p(f(v))-\lceil N(P)/2\rceil-j.
\]
Since \(N(\nu)=N(P)+2j\), the potential identity gives the first
congruence.  Applying it at the entry event gives the second.
For a finite \(\omega\)-tail the coefficient decreases by one at every
event, giving the exact formulas directly.
\end{proof}

For use in arithmetic elimination, the complete tables are as follows.
Here \(a\ge0\), and every displayed value residue has modulus \(2m\)
except in the last row, whose modulus is \(m\).
\begin{center}\small
\begin{tabular}{@{}lll@{}}\toprule
Field & \(k\)-event interval & Other proper interval\\\midrule
\(\tau(t)\pmod m\) & \(k\) & \(p(t)-c_\nu\)\\
\(E(t)\pmod{2m}\) & \(2p(t)-2k\) & \(2c_\nu\)\\
\(f^a(E(t))\pmod{2m}\) & \(2p(t)-2(k+a)\) & \(2c_\nu-2a\)\\
\(f(E(t))/2\pmod m\) & \(p(t)-k-1\) & \(c_\nu-1\)\\\bottomrule
\end{tabular}
\end{center}
For a bare iterate \(f^a(t)\) with \(a\ge1\), the residue modulo
\(2m\) is \(2p(t)-2a\).  The identity field remains \(t\), not
\(2p(t)\); in particular its odd bit must be retained.

The finite prefix cannot be discarded.  If \(b\ge3\), then at
\(t=3\) the preserved control gives
\(D_{\omega^2}(3)=G(3)/2+1\equiv0\pmod3\), whereas
\(p(3)=1\).  The value \(G(3)\) is a power of four, which proves
the displayed congruence without evaluating it.  Small inputs will be
handled by the exact evaluator, not by extending the tail tables.
